\section{Limitations}
\label{sec:limitations}

\subsection{Internal validity}
\label{sec:limits-internal}

Every number comes from one generator with one volatility law, a log-AR(1) of persistence 0.98 that is known at the start of the day, and one pricing noise of 10\% of the factor profit and loss.
The stale-volatility deficiency is the easiest form of a scale-timing defect to detect: the lag is exact, the volatility is persistent enough that a lag of 20 days leaves the correlation of the rolling means at 0.69, and the generator's own risk model is otherwise correct.
The thresholds are specific to that generator, and the alternative nulls of Section~\ref{sec:results-altnull} show the consequence: the size is 4.0 to 6.2\% on two of the four alternatives that change the volatility law and 99.0\% on one that changes the pricing noise.
The ablation of Section~\ref{sec:ablation} was computed on the committed seeds that the main results use, so a choice made from it, such as a shorter window, would be selected on the evaluation data.
The three comparators that the engine does not carry were written for this evaluation; they agree with independent implementations to a relative difference below $10^{-9}$ on a desk-year, which checks their arithmetic and not whether the published tests are best applied this way (the Ljung--Box lag of 10 and the squared residual are choices of ours, not of the cited authors).
The Monte Carlo spread of the critical values, a standard deviation of 0.10 percentage points at a size of 1\% and 0.058 at 0.333\% (Proposition~\ref{prop:size}), is small against the effects reported but not against the difference between two sizes of 1.0 and 1.2\%, which the tables do not claim to resolve.

\subsection{External validity}
\label{sec:limits-external}

The desks are synthetic: two factors, one volatility process, one deficiency at a time and a window of 250 days.
No real desk is evaluated, and the results say nothing about the correlations, the distances or the pricing noise of the desks of a bank, nor about how often a real risk model has a stale volatility input.
A real desk has many factors, jumps in the profit and loss, a pricing noise that depends on the instruments and the market regime, several simultaneous deficiencies, and a back-history that may be shorter than the null the calibration needs.
The calibration needs a simulator of the desk's own null.
A bank that has one can use the complement as specified; a bank that does not would have to build the null from a model of its desk or from historical windows in which the model is known to be sound, and neither route was tested here.
The proposition on the size requires the desk-year to be exchangeable with the null draws, and the experiments show what the size becomes when it is not.

\subsection{Construct validity and scope}
\label{sec:limits-construct}

The detection rate of a planted deficiency is a measure of the statistic's power against that deficiency in that generator.
It is not a measure of usefulness for the supervisory decision, in which the costs of a false alarm and of a miss are not those of the experiments.
The lag of 20 days is one way of getting the scale wrong; a misparametrised decay of an exponentially weighted volatility, a volatility that omits a regime, and a volatility estimated on a window of the wrong length are others, and they were not planted.
The complement has no regulatory standing: it changes neither the regulatory zone nor the thresholds of the standard \citep{bcbs2019}, which the engine returns unchanged beside it, and the manuscript makes no claim about how a supervisor should use it.
The comparison with published methods covers the nine detectors of Section~\ref{sec:design-comparators}.
The review of related work is limited to the sources cited, and the absence of a published statistic of the same form as the complement is a statement about that search and not a proof.
